%! AP Statistics — Unit 4, Lesson 4.1 — Exit Ticket KEY
\documentclass[10pt]{article}
\usepackage{apstats-article}
\usepackage{apstats-key}

\begin{document}

\pageheader{Unit 4, Lesson 4.1}{Exit Ticket --- KEY}
\namedateperiod

\vspace{0.14in}

\begin{scenariobox}[Context]{sky}
\small
A student rolls a standard six-sided die 20 times and records:\\[4pt]
3,\ 3,\ 3,\ 3,\ 5,\ 2,\ 1,\ 6,\ 3,\ 4,\ 5,\ 2,\ 3,\ 1,\ 6,\ 4,\ 4,\ 4,\ 4,\ 4
\end{scenariobox}

\vspace{0.1in}

\begin{practicebox}
\small
\begin{enumerate}[leftmargin=1.5em, itemsep=8pt]
  \item
        \begin{enumerate}[label=(\alph*), itemsep=4pt]
          \item \ans{3 appears 6 times (30\%), far above the expected $\frac{1}{6} \approx 17\%$; there is a streak of four 3s at the start.}
          \item \ans{4 appears 5 times (25\%) with a streak of five consecutive 4s at the end.}
        \end{enumerate}

  \item \ans{Is the die biased toward 3 and 4, or could the pattern of repeated 3s and 4s occur by chance with a fair die?}

  \item \ans{With only 20 rolls, random variation can easily produce clusters of the same face. A fair die \emph{will} produce runs; that is a natural feature of random processes. To determine whether the die is actually unfair, we would need many more rolls and a formal probability calculation to determine how unlikely these counts are if the die were fair.}
\end{enumerate}
\end{practicebox}

\end{document}
